\begin{thebibliography}{9}
\raggedright
\phantomsection
\addcontentsline{toc}{section}{Referencias}

\bibitem{edaa-sparse-table}
LELE. \textit{Generic Sparse Table} [Código fuente].
Repositorio del curso EDAA.
\url{https://github.com/jfuentess/edaa/blob/550e2199f82e74e38d253eba2e64f8c72af38296/rmq/sparse_table.cpp}.

\bibitem{edaa-segment-tree}
LELE. \textit{Generic Segment Tree for RMQs} [Código fuente].
Repositorio del curso EDAA.
\url{https://github.com/jfuentess/edaa/blob/550e2199f82e74e38d253eba2e64f8c72af38296/rmq/segment_tree.cpp}.

\bibitem{cp-sparse-table}
Colaboradores de CP-Algorithms.
\textit{Sparse Table}.
\url{https://cp-algorithms.com/data_structures/sparse-table.html}.

\bibitem{cp-segment-tree}
Colaboradores de CP-Algorithms.
\textit{Segment Tree}.
\url{https://cp-algorithms.com/data_structures/segment_tree.html}.

\bibitem{atcoder-segtree}
AtCoder Library.
\textit{Segtree}: documentación de árboles de segmentos sobre monoides.
\url{https://atcoder.github.io/ac-library/production/document_en/segtree.html}.

\bibitem{cpp-bit-width}
Cppreference.
\textit{std::bit\_width}.
\url{https://en.cppreference.com/w/cpp/numeric/bit_width.html}.

\bibitem{cp-euclid}
Colaboradores de CP-Algorithms.
\textit{Euclidean algorithm for computing the greatest common divisor}.
\url{https://cp-algorithms.com/algebra/euclid-algorithm.html}.

\bibitem{cpp-lcm}
Cppreference.
\textit{std::lcm}.
\url{https://en.cppreference.com/w/cpp/numeric/lcm.html}.

\bibitem{lovera-uhr}
Lovera, L. (LELE). (s.\,f.).
\textit{uhr: generic time performance tester} [Código fuente].
Plantilla de medición de la ayudantía, disponible en
\texttt{utils/uhr.cpp} del repositorio de estos boletines.
Referencia del autor utilizada en los boletines anteriores:
\url{https://github.com/leonardlover}.

\par\medskip
\noindent\textbf{Declaración de uso de IA generativa.}
Se utilizaron herramientas de inteligencia artificial generativa como apoyo para las implementaciones de los experimentos y para revisar la redacción técnica de este informe. Codex se utilizó como agente prioritario durante este proceso. El autor es responsable de las decisiones finales, la verificación de los resultados y el contenido presentado.

\end{thebibliography}
